% ============================================================
%  PREAMBOLO DELLA DISPENSA DI METODI STATISTICI PER L'INFORMAZIONE
% ============================================================
% Oltre ai riquadri della classe (\dfn, \info, \warning, \error, \gbox) la
% dispensa usa cinque ambienti pensati per la matematica:
%   \teorema[etichetta]{Titolo}{enunciato}   teorema numerato (blu)
%   dimostrazione                            dimostrazione passo per passo, con \passo{...}
%   richiamo{Titolo}                         richiamo di matematica: derivate, integrali, limiti... (viola)
%   formula{Titolo}                          formula da ricordare, in evidenza (verde petrolio)
%   esempio{Titolo} / esercizio{Titolo}      esempi ed esercizi numerati per capitolo
% Le dimostrazioni e i richiami sono ambienti (non comandi): possono contenere
% figure TikZ, elenchi e formule su piu' righe, e si spezzano fra due pagine.

% ------------------------------------------------------------
% Pacchetti
% ------------------------------------------------------------
\usepackage{pgfplots}
\pgfplotsset{compat=1.18}
\usepgfplotslibrary{fillbetween, groupplots}
\usepackage{booktabs}
\usepackage{multirow}
\usepackage{colortbl}
\usepackage{needspace}
\usepackage{cancel}
\usepackage{icomma}                 % in matematica la virgola decimale non lascia spazi: 0,5
% Didascalie come nella dispensa di Tecnologie Web (il pacchetto evita anche le ancore
% duplicate che le figure [H] della classe producono nel PDF)
\usepackage[font=small,labelfont={bf,color=primary!80!black},
            textfont=it,skip=6pt]{caption}
\usetikzlibrary{calc, arrows.meta, patterns, positioning, shapes.geometric, shapes.misc,
                decorations.pathreplacing, decorations.pathmorphing, shadows, fit, backgrounds}
\tcbuselibrary{raster}

\newcolumntype{L}{>{\raggedright\arraybackslash}X}
\newcolumntype{P}[1]{>{\raggedright\arraybackslash}p{#1}}
\newcolumntype{C}[1]{>{\centering\arraybackslash}p{#1}}

% ------------------------------------------------------------
% Colori
% ------------------------------------------------------------
\definecolor{mspiteo}{HTML}{1F5FA8}     % teoremi
\definecolor{mspidim}{HTML}{4A5568}     % dimostrazioni
\definecolor{mspiric}{HTML}{7B3FA0}     % richiami di matematica
\definecolor{mspiform}{HTML}{0B7A5A}    % formule da ricordare
\definecolor{mspiA}{HTML}{2B6CB0}       % evento A, prima curva
\definecolor{mspiB}{HTML}{C05621}       % evento B, seconda curva
\definecolor{mspiC}{HTML}{2F855A}       % terzo evento, terza curva
\definecolor{mspigray}{HTML}{718096}
\definecolor{mspilight}{HTML}{EDF2F7}

% ------------------------------------------------------------
% Titoli e riquadri che non si spezzano
% ------------------------------------------------------------
% Pagine di apertura delle parti nello stile del documento (come nella dispensa di RPD)
\titleformat{\part}[display]
  {\centering\color{primary}}
  {\fontsize{20}{24}\selectfont\Gobold Parte \thepart}
  {18pt}
  {\color{primary!80!black}\fontsize{32}{40}\selectfont\Gobold #1}
  [\vspace{10pt}{\color{primary!40}\rule{0.35\linewidth}{2pt}}]
\assignpagestyle{\part}{empty}

% Indice: i numeri di pagina a tre cifre non entrano nella casella standard di 1,55em e la
% classe (\overfullrule) segnerebbe ogni riga con un quadratino nero; la casella si allarga.
\makeatletter
\renewcommand{\@pnumwidth}{2.2em}
\makeatother

\let\oldsection\section
\renewcommand{\section}{\Needspace{16\baselineskip}\oldsection}
\let\oldsubsection\subsection
\renewcommand{\subsection}{\Needspace{10\baselineskip}\oldsubsection}

% \mspimisura{contenuto}: compone il contenuto in una scatola invisibile solo per
% misurarne l'altezza; se sta in una pagina chiede quello spazio con \Needspace, così il
% riquadro vero, composto subito dopo, passa intero alla pagina successiva invece di
% spezzarsi. Se è più alto di una pagina si spezza normalmente. (Dentro i riquadri
% misurati non si mettono \label: verrebbero definite due volte.)
\newbox\mspimisurabox
\newcommand{\mspimisura}[1]{%
  \par
  \setbox\mspimisurabox\vbox{#1}%
  \ifdim\dimexpr\ht\mspimisurabox+\dp\mspimisurabox\relax<\textheight
    % +1cm: lo spazio che il riquadro lascia sopra di sé (before skip) non è nella misura
    \Needspace{\dimexpr\ht\mspimisurabox+\dp\mspimisurabox+1cm\relax}%
  \fi}

% Le definizioni non si spezzano mai; gli altri riquadri passano interi alla pagina dopo.
% Il titolo va fra graffe: una virgola nel titolo non viene presa per un separatore di opzioni.
\renewcommand{\dfn}[2]{\Needspace{6\baselineskip}\begin{Definition*}[colbacktitle=red!75!black, breakable=false]{#1}{}#2\end{Definition*}}
\renewcommand{\info}[2]{\mspimisura{\begin{infobox}{{#1}}#2\end{infobox}}\begin{infobox}{{#1}}#2\end{infobox}}
\renewcommand{\warning}[2]{\mspimisura{\begin{warningbox}{{#1}}#2\end{warningbox}}\begin{warningbox}{{#1}}#2\end{warningbox}}
\renewcommand{\error}[2]{\mspimisura{\begin{errorbox}{{#1}}#2\end{errorbox}}\begin{errorbox}{{#1}}#2\end{errorbox}}
\let\oldgbox\gbox
\renewcommand{\gbox}[3]{\mspimisura{\oldgbox{#1}{#2}{#3}}\oldgbox{#1}{#2}{#3}}

% ------------------------------------------------------------
% Teoremi
% ------------------------------------------------------------
% Stessa forma delle definizioni della classe, in blu e numerati per capitolo:
% «Teorema 3.2 --- Teorema di Bayes».
\newtcbtheorem[number within=chapter]{teoremabox}{Teorema}{
  enhanced, breakable=false,
  before skip=5mm, after skip=5mm,
  colback=mspiteo!4, colframe=mspiteo!85!black, boxrule=0.5mm,
  colbacktitle=mspiteo!85!black, coltitle=white, fonttitle=\bfseries,
  separator sign={\ ---}, description delimiters none,
  attach boxed title to top left={xshift=1cm,yshift*=1mm-\tcboxedtitleheight},
  drop fuzzy shadow,
  boxed title style={
    frame code={
      \path[fill=tcbcolback]
      ([yshift=-1mm,xshift=-1mm]frame.north west)
      arc[start angle=0,end angle=180,radius=1mm]
      ([yshift=-1mm,xshift=1mm]frame.north east)
      arc[start angle=180,end angle=0,radius=1mm];
      \path[left color=tcbcolback!60!black,right color=tcbcolback!60!black,
        middle color=tcbcolback!80!black]
      ([xshift=-2mm]frame.north west) -- ([xshift=2mm]frame.north east)
           [rounded corners=1mm]-- ([xshift=1mm,yshift=-1mm]frame.north east)
           -- (frame.south east) -- (frame.south west)
           -- ([xshift=-1mm,yshift=-1mm]frame.north west)
           [sharp corners]-- cycle;
    },
    interior engine=empty,
  },
}{teo}
% \teorema[etichetta]{Titolo}{enunciato}: con l'etichetta si cita come \ref{teo:etichetta}
\NewDocumentCommand{\teorema}{o m m}{%
  \Needspace{8\baselineskip}%
  \IfNoValueTF{#1}{\begin{teoremabox}{#2}{}#3\end{teoremabox}}%
                  {\begin{teoremabox}{#2}{#1}#3\end{teoremabox}}}

% ------------------------------------------------------------
% Dimostrazioni
% ------------------------------------------------------------
% Riquadro con la barra a sinistra, come le note della classe, in grigio ardesia.
% Finisce con il quadratino nero che segna «come volevasi dimostrare».
% Dentro, \passo{Titolo} numera i passaggi: «Passo 1 --- Titolo».
\newcounter{passo}
\newcommand{\passo}[1]{%
  \par\smallskip\stepcounter{passo}%
  \noindent{\color{mspidim}\bfseries Passo \thepasso\ifx\relax#1\relax\else\ --- #1\fi.}\ }
\newtcolorbox{dimostrazionebox}[1][]{%
  title={Dimostrazione\ifx\relax#1\relax\else\ --- #1\fi}, fonttitle=\bfseries,
  drop fuzzy shadow,
  colback=mspilight!60!white, enhanced, breakable,
  attach boxed title to top left={xshift=-5mm},
  boxrule=0pt, before skip=0.8cm, after skip=0.8cm,
  toprule=0pt, bottomrule=0pt, rightrule=0pt, leftrule=4pt,
  skin=enhancedlast jigsaw, colframe=mspidim, colbacktitle=mspidim!85!white,
  before upper={\setcounter{passo}{0}},
  after upper={\par\hfill{\color{mspidim}$\blacksquare$}},
  boxed title style={
    frame code={
      \fill[mspidim!85!white](frame.south west)--(frame.north west)--(frame.north east)--([xshift=3mm]frame.east)--(frame.south east)--cycle;
      \draw[line width=1mm,mspidim!85!white]([xshift=2mm]frame.north east)--([xshift=5mm]frame.east)--([xshift=2mm]frame.south east);
      \draw[line width=1mm,mspidim!85!white]([xshift=5mm]frame.north east)--([xshift=8mm]frame.east)--([xshift=5mm]frame.south east);
      \fill[mspidim!40](frame.south west)--+(4mm,-2mm)--+(4mm,2mm)--cycle;
    }
  },
}
\NewDocumentEnvironment{dimostrazione}{O{} +b}
  {\mspimisura{\begin{dimostrazionebox}[#1]#2\end{dimostrazionebox}}%
   \begin{dimostrazionebox}[#1]#2\end{dimostrazionebox}}{}

% ------------------------------------------------------------
% Richiami di matematica
% ------------------------------------------------------------
% Le spiegazioni «da zero» di limiti, derivate, integrali, sommatorie...: stessa
% forma delle note della classe, in viola, con l'icona del libro.
\newtcolorbox{richiamobox}[1]{%
  title={\ibook\ \ Richiamo di matematica: #1}, fonttitle=\bfseries,
  drop fuzzy shadow,
  colback=mspiric!5, enhanced, breakable,
  attach boxed title to top left={xshift=-5mm},
  boxrule=0pt, before skip=0.8cm, after skip=0.8cm,
  toprule=0pt, bottomrule=0pt, rightrule=0pt, leftrule=4pt,
  skin=enhancedlast jigsaw, colframe=mspiric, colbacktitle=mspiric!80!white,
  boxed title style={
    frame code={
      \fill[mspiric!80!white](frame.south west)--(frame.north west)--(frame.north east)--([xshift=3mm]frame.east)--(frame.south east)--cycle;
      \draw[line width=1mm,mspiric!80!white]([xshift=2mm]frame.north east)--([xshift=5mm]frame.east)--([xshift=2mm]frame.south east);
      \draw[line width=1mm,mspiric!80!white]([xshift=5mm]frame.north east)--([xshift=8mm]frame.east)--([xshift=5mm]frame.south east);
      \fill[mspiric!40](frame.south west)--+(4mm,-2mm)--+(4mm,2mm)--cycle;
    }
  },
}
\NewDocumentEnvironment{richiamo}{m +b}
  {\mspimisura{\begin{richiamobox}{#1}#2\end{richiamobox}}%
   \begin{richiamobox}{#1}#2\end{richiamobox}}{}

% ------------------------------------------------------------
% Formule da ricordare
% ------------------------------------------------------------
\newtcolorbox{formula}[1][]{%
  enhanced, breakable=false, center, width=0.9\linewidth,
  colback=mspiform!5, colframe=mspiform!80!black, boxrule=0.6pt, arc=2mm,
  before skip=4mm, after skip=4mm, left=6pt, right=6pt, top=4pt, bottom=4pt,
  fonttitle=\small\bfseries, coltitle=mspiform!80!black, colbacktitle=mspiform!5,
  title={\ifx\relax#1\relax Da ricordare\else Da ricordare: #1\fi},
  attach title to upper={\par\smallskip}, drop fuzzy shadow,
}

% ------------------------------------------------------------
% Esempi ed esercizi
% ------------------------------------------------------------
\setcounter{secnumdepth}{0}
\makeatletter
\@addtoreset{esempio}{chapter}
\@addtoreset{esercizio}{chapter}
\makeatother
% Composti in una scatola: se stanno in una pagina non si spezzano mai (passano interi
% alla pagina successiva), altrimenti la scatola viene «svuotata» e si spezzano normalmente.
\newbox\mspiesbox
\newcommand{\mspiemetti}{%
  \ifdim\dimexpr\ht\mspiesbox+\dp\mspiesbox\relax>\textheight
    \unvbox\mspiesbox
  \else
    \box\mspiesbox
  \fi}
\renewenvironment{esempio}[1]{%
  \par\addvspace{0.8em}%
  \setbox\mspiesbox\vbox\bgroup
  \refstepcounter{esempio}%
  \noindent\rule{\linewidth}{0.5pt}\par\nopagebreak
  \vspace{0.15em}%
  \noindent{\color{primary!80!black}\textbf{Esempio \thechapter.\theesempio}}%
  \ifx\relax#1\relax\else{\color{primary!80!black}\textbf{ --- #1}}\fi\par\nopagebreak
  \vspace{0.3em}%
}{%
  \par\vspace{0.3em}%
  \noindent\rule{\linewidth}{0.5pt}\par
  \egroup
  \mspiemetti
  \par\vspace{0.8em}%
}
\renewenvironment{esercizio}[1]{%
  \par\addvspace{0.8em}%
  \setbox\mspiesbox\vbox\bgroup
  \refstepcounter{esercizio}%
  \noindent\rule{\linewidth}{1.2pt}\vspace{1pt}\par
  \noindent\rule{\linewidth}{0.4pt}\par
  \vspace{0.15em}%
  \noindent\textbf{Esercizio \thechapter.\theesercizio}%
  \ifx\relax#1\relax\else\textbf{ --- #1}\fi\par
  \vspace{0.3em}%
}{%
  \par\vspace{0.3em}%
  \noindent\rule{\linewidth}{0.4pt}\vspace{1pt}\par
  \noindent\rule{\linewidth}{1.2pt}\par
  \egroup
  \mspiemetti
  \par\vspace{0.8em}%
}
% Intestazione dello svolgimento di un esercizio
\newcommand{\svolgimento}{\par\smallskip\noindent{\sffamily\bfseries\color{primary!80!black}Svolgimento}\par\nobreak\smallskip}

% ------------------------------------------------------------
% Notazione
% ------------------------------------------------------------
\providecommand{\term}[1]{\textbf{#1}}                % termine introdotto
\providecommand{\R}{\mathbb{R}}
\providecommand{\N}{\mathbb{N}}
\providecommand{\Z}{\mathbb{Z}}
\providecommand{\dd}{\,\mathrm{d}}                         % il «d» di dx negli integrali
\providecommand{\abs}[1]{\left\lvert #1 \right\rvert}
% La classe ridefinisce \abs con \@ifstar, che non è robusto: dentro una didascalia
% (\caption) si rompe. Lo si ridefinisce robusto, con lo stesso comportamento
% (\abs adatta le barre all'altezza del contenuto, \abs* no).
\makeatletter
\expandafter\let\csname mspi@absimpl\expandafter\endcsname\csname abs \endcsname
\DeclareRobustCommand{\abs}{\@ifstar{\mspi@absimpl}{\mspi@absimpl*}}
\makeatother
\providecommand{\card}[1]{\left\lvert #1 \right\rvert}     % numero di elementi di un insieme
\providecommand{\compl}[1]{\overline{#1}}                  % complementare di un evento
% Come nelle slide: vettori sottolineati una volta, matrici due volte. I \mkern lasciano un
% piccolo stacco fra le sottolineature di due simboli vicini (in A x non diventano una riga sola).
\newcommand{\vett}[1]{{\mkern1mu\underline{#1}\mkern1mu}}
\newcommand{\matr}[1]{{\mkern1mu\underline{\underline{#1}}\mkern1mu}}

% ------------------------------------------------------------
% Grafici
% ------------------------------------------------------------
% Assi «da quaderno»: frecce, origine al centro, niente cornice.
\pgfplotsset{
  quaderno/.style={
    scaled ticks=false,
    xticklabel style={/pgf/number format/fixed, /pgf/number format/use comma},
    yticklabel style={/pgf/number format/fixed, /pgf/number format/use comma},
    axis lines=middle, axis line style={-{Stealth[length=2mm]}, line width=0.6pt},
    xlabel style={anchor=north west}, ylabel style={anchor=south east},
    tick label style={font=\scriptsize}, label style={font=\small},
    grid=major, grid style={mspigray!20},
    every axis plot/.append style={line width=1.1pt},
    clip=false,
  },
  % Distribuzioni discrete: bastoncini con il pallino in cima
  bastoncini/.style={ycomb, mark=*, mark size=1.6pt, line width=1.4pt},
  % Realizzazioni di processi: bastoncini che partono dalla linea dello zero, numeri del tempo in basso
  realizzazione/.style={
    axis x line=bottom, axis y line=left, axis line style={-}, tick label style={font=\scriptsize},
    label style={font=\small}, extra y ticks={0},
    extra y tick style={grid=major, major grid style={mspigray, line width=0.5pt}},
    every axis plot/.append style={ycomb, mark=*, mark size=1.3pt, line width=0.9pt}, clip=false,
  },
}

% ------------------------------------------------------------
% Disegni ricorrenti
% ------------------------------------------------------------
\colorlet{evid}{yellow!55}       % riempimento delle regioni evidenziate

% Facce del dado: \dado{n}, n da 1 a 6; \dado[lato]{n} cambia la dimensione
\NewDocumentCommand{\dado}{O{0.42cm} m}{%
  \tikz[baseline=-0.65ex, x=#1, y=#1]{%
    \draw[rounded corners=0.6mm, fill=white, line width=0.5pt] (-0.5,-0.5) rectangle (0.5,0.5);
    \ifcase#2\or
      \fill (0,0) circle (0.1);
    \or
      \fill (-0.24,0.24) circle (0.1) (0.24,-0.24) circle (0.1);
    \or
      \fill (-0.25,0.25) circle (0.1) (0,0) circle (0.1) (0.25,-0.25) circle (0.1);
    \or
      \fill (-0.24,0.24) circle (0.1) (0.24,0.24) circle (0.1)
            (-0.24,-0.24) circle (0.1) (0.24,-0.24) circle (0.1);
    \or
      \fill (-0.25,0.25) circle (0.1) (0.25,0.25) circle (0.1) (0,0) circle (0.1)
            (-0.25,-0.25) circle (0.1) (0.25,-0.25) circle (0.1);
    \or
      \fill (-0.24,0.27) circle (0.1) (-0.24,0) circle (0.1) (-0.24,-0.27) circle (0.1)
            (0.24,0.27) circle (0.1) (0.24,0) circle (0.1) (0.24,-0.27) circle (0.1);
    \fi}}

% Una moneta: \moneta{T} o \moneta{C}
\newcommand{\moneta}[1]{%
  \tikz[baseline=-0.65ex]{\draw[fill=yellow!35, draw=yellow!40!black, line width=0.5pt] (0,0) circle (0.2cm);
    \node[font=\scriptsize\bfseries, yellow!25!black, inner sep=0pt] at (0,0) {#1};}}

% Diagrammi di Venn con due eventi A e B dentro Omega.
% \venn{riempimenti}{didascalia}: i riempimenti si scrivono con le forme qui sotto, per esempio
%   \fill[evid] \vennA;                                         (l'evento A)
%   \begin{scope}\clip \vennA; \fill[evid] \vennB; \end{scope}   (l'intersezione)
\newcommand{\vennOmega}{(-1.5,-1.0) rectangle (1.5,1.0)}
\newcommand{\vennA}{(-0.4,0) circle (0.68)}
\newcommand{\vennB}{(0.4,0) circle (0.68)}
\NewDocumentCommand{\venn}{m m}{%
  \begin{tikzpicture}[baseline=(current bounding box.north)]
    #1
    \draw[mspigray, line width=0.5pt, rounded corners=2pt] \vennOmega;
    \draw[mspiA!85!black, line width=0.8pt] \vennA;
    \draw[mspiB!85!black, line width=0.8pt] \vennB;
    \node[mspiA!85!black, font=\small] at (-1.12,0.68) {$A$};
    \node[mspiB!85!black, font=\small] at (1.12,0.68) {$B$};
    \node[mspigray, font=\footnotesize, anchor=south east, inner sep=2pt] at (1.5,-1.0) {$\Omega$};
    \node[font=\small, anchor=north, text width=3.2cm, align=center] at (0,-1.08) {#2};
  \end{tikzpicture}}

% Lo stesso con tre eventi A, B, C
\newcommand{\vennOmegaT}{(-1.6,-1.45) rectangle (1.6,1.25)}
\newcommand{\vennTA}{(-0.4,0.25) circle (0.68)}
\newcommand{\vennTB}{(0.4,0.25) circle (0.68)}
\newcommand{\vennTC}{(0,-0.42) circle (0.68)}
\NewDocumentCommand{\venntre}{m m}{%
  \begin{tikzpicture}[baseline=(current bounding box.north)]
    #1
    \draw[mspigray, line width=0.5pt, rounded corners=2pt] \vennOmegaT;
    \draw[mspiA!85!black, line width=0.8pt] \vennTA;
    \draw[mspiB!85!black, line width=0.8pt] \vennTB;
    \draw[mspiC!75!black, line width=0.8pt] \vennTC;
    \node[mspiA!85!black, font=\small] at (-1.18,0.9) {$A$};
    \node[mspiB!85!black, font=\small] at (1.18,0.9) {$B$};
    \node[mspiC!75!black, font=\small] at (0.95,-1.05) {$C$};
    \node[mspigray, font=\footnotesize, anchor=south west, inner sep=2pt] at (-1.6,-1.45) {$\Omega$};
    \node[font=\small, anchor=north, text width=3.4cm, align=center] at (0,-1.53) {#2};
  \end{tikzpicture}}

% ------------------------------------------------------------
% Pagine
% ------------------------------------------------------------
\raggedbottom
% Niente righe isolate in cima o in fondo alla pagina. Le penalità vanno impostate a
% documento iniziato: all'inizio del documento vengono riportate a 3000.
\makeatletter
\AddToHook{begindocument/end}{%
  \widowpenalty=10000 \clubpenalty=10000 \@clubpenalty=10000 \brokenpenalty=10000}
\makeatother
\emergencystretch=2.5em
